\section*{References}
{\small\setlength{\parindent}{-0.25in}\setlength{\leftskip}{0.25in}\setlength{\parskip}{3pt}
Austin, P. C., \& Stuart, E. A. (2015). Moving towards best practice when using inverse probability of treatment weighting (IPTW) using the propensity score to estimate causal treatment effects in observational studies. Statistics in Medicine, 34(28), 3661-3679. \url{https://doi.org/10.1002/sim.6607}

Bang, H., \& Robins, J. M. (2005). Doubly robust estimation in missing data and causal inference models. Biometrics, 61(4), 962-973. \url{https://doi.org/10.1111/j.1541-0420.2005.00377.x}

Chernozhukov, V., Chetverikov, D., Demirer, M., Duflo, E., Hansen, C., Newey, W., \& Robins, J. (2018). Double/debiased machine learning for treatment and structural parameters. The Econometrics Journal, 21(1), C1-C68. \url{https://doi.org/10.1111/ectj.12097}

Dietz, P. M., Homa, D., England, L. J., Burley, K., Tong, V. T., Dube, S. R., \& Bernert, J. T. (2011). Estimates of nondisclosure of cigarette smoking among pregnant and nonpregnant women of reproductive age in the United States. American Journal of Epidemiology, 173(3), 355-359. \url{https://doi.org/10.1093/aje/kwq381}

Funk, M. J., Westreich, D., Wiesen, C., Stürmer, T., Brookhart, M. A., \& Davidian, M. (2011). Doubly robust estimation of causal effects. American Journal of Epidemiology, 173(7), 761-767. \url{https://doi.org/10.1093/aje/kwq439}

Howland, R. E., Mulready-Ward, C., Madsen, A. M., Sackoff, J., Nyland-Funke, M., Bombard, J. M., \& Tong, V. T. (2015). Reliability of reported maternal smoking: Comparing the birth certificate to maternal worksheets and prenatal and hospital medical records, New York City and Vermont, 2009. Maternal and Child Health Journal, 19(9), 1916-1924. \url{https://doi.org/10.1007/s10995-015-1722-1}

Joyce, T., Racine, A., \& Yunzal-Butler, C. (2008). Reassessing the WIC effect: Evidence from the Pregnancy Nutrition Surveillance System. Journal of Policy Analysis and Management, 27(2), 277-303.

Jurek, A. M., Greenland, S., Maldonado, G., \& Church, T. R. (2005). Proper interpretation of non-differential misclassification effects: Expectations vs observations. International Journal of Epidemiology, 34(3), 680-687. \url{https://doi.org/10.1093/ije/dyi060}

Kondracki, A. J. (2020). Low birthweight in term singletons mediates the association between maternal smoking intensity exposure status and immediate neonatal intensive care unit admission: The E-value assessment. BMC Pregnancy and Childbirth. \url{https://doi.org/10.1186/s12884-020-02981-1}

Leung, M., Kioumourtzoglou, M.-A., Raz, R., \& Weisskopf, M. G. (2021). Bias due to selection on live births in studies of environmental exposures during pregnancy: A simulation study. Environmental Health Perspectives, 129(4), 047001. \url{https://doi.org/10.1289/EHP7961}

Li, F., Morgan, K. L., \& Zaslavsky, A. M. (2018). Balancing covariates via propensity score weighting. Journal of the American Statistical Association, 113(521), 390-400. \url{https://doi.org/10.1080/01621459.2016.1260466}

Liu, B., Xu, G., Sun, Y., Qiu, X., Ryckman, K. K., Yu, Y., Snetselaar, L. G., \& Bao, W. (2020). Maternal cigarette smoking before and during pregnancy and the risk of preterm birth: A dose-response analysis of 25 million mother-infant pairs. PLOS Medicine, 17(8), e1003158. \url{https://doi.org/10.1371/journal.pmed.1003158}

Martin, J. A., Hamilton, B. E., \& Osterman, M. J. K. (2024). Births in the United States, 2023 (NCHS Data Brief No. 507). National Center for Health Statistics. \url{https://doi.org/10.15620/cdc/158789}

Martin, J. A., Osterman, M. J. K., \& Driscoll, A. K. (2023). Declines in cigarette smoking during pregnancy in the United States, 2016-2021 (NCHS Data Brief No. 458). National Center for Health Statistics.

Martin, J. A., Osterman, M. J. K., Kirmeyer, S. E., \& Gregory, E. C. W. (2015). Measuring gestational age in vital statistics data: Transitioning to the obstetric estimate. National Vital Statistics Reports, 64(5).

National Center for Health Statistics. (2024). User guide to the 2023 natality public use file. \url{https://ftp.cdc.gov/pub/Health_Statistics/NCHS/Dataset_Documentation/DVS/natality/UserGuide2023.pdf}

Patra, J., Bakker, R., Irving, H., Jaddoe, V. W. V., Malini, S., \& Rehm, J. (2011). Dose-response relationship between alcohol consumption before and during pregnancy and the risks of low birthweight, preterm birth and small for gestational age (SGA): A systematic review and meta-analyses. BJOG, 118(12), 1411-1421. \url{https://doi.org/10.1111/j.1471-0528.2011.03050.x}

Pearl, J. (1995). Causal diagrams for empirical research. Biometrika, 82(4), 669-688. \url{https://doi.org/10.1093/biomet/82.4.669}

Regan, A. K., Bombard, J. M., O'Hegarty, M. M., Smith, R. A., \& Tong, V. T. (2021). Adverse birth outcomes associated with prepregnancy and prenatal electronic cigarette use. Obstetrics \& Gynecology, 138(1), 85-94. \url{https://doi.org/10.1097/AOG.0000000000004432}

Searles Nielsen, S., Dills, R. L., Glass, M., \& Mueller, B. A. (2014). Accuracy of prenatal smoking data from Washington State birth certificates in a population-based sample with cotinine measurements. Annals of Epidemiology, 24(3). \url{https://doi.org/10.1016/j.annepidem.2013.12.008}

Shin, D., \& Song, W. O. (2019). Influence of the adequacy of the Prenatal Care Utilization Index on small-for-gestational-age infants and preterm births in the United States. Journal of Clinical Medicine, 8(6), 838. \url{https://doi.org/10.3390/jcm8060838}

Stock, S. J., \& Bauld, L. (2020). Maternal smoking and preterm birth: An unresolved health challenge. PLOS Medicine. \url{https://doi.org/10.1371/journal.pmed.1003386}

Tennant, P. W. G., Murray, E. J., Arnold, K. F., Berrie, L., Fox, M. P., Gadd, S. C., ... Ellison, G. T. H. (2021). Use of directed acyclic graphs (DAGs) to identify confounders in applied health research: Review and recommendations. International Journal of Epidemiology, 50(2), 620-632. \url{https://doi.org/10.1093/ije/dyaa213}

VanderWeele, T. J., \& Ding, P. (2017). Sensitivity analysis in observational research: Introducing the E-value. Annals of Internal Medicine, 167(4), 268-274. \url{https://doi.org/10.7326/M16-2607}

Zivich, P. N., \& Breskin, A. (2021). Machine learning for causal inference: On the use of cross-fit estimators. Epidemiology, 32(3), 393-401. \url{https://doi.org/10.1097/EDE.0000000000001332}
}
